% LEGACY, UNCOMPILED DRAFT FRAGMENT. The authoritative manuscript is ../main.tex.
\section{Introduction}
\label{sec:introduction}

Brain tumors alter local anatomy and violate the healthy-tissue assumptions of many
registration, parcellation, and atlas-based neuroimaging pipelines. Counterfactual
inpainting replaces a pathological region with plausible tumor-free tissue so that
such tools can operate on a completed image. The BraTS local-synthesis
task~\cite{kofler2023inpainting} formalizes this problem: given a voided T1n
(\emph{T1-weighted native}) volume and its inpainting mask, synthesize the missing
anatomy while leaving all observed voxels unchanged.

Diffusion models~\cite{ho2020ddpm} provide effective conditional
generators~\cite{lugmayr2022repaint,saharia2022palette}. Direct full-resolution 3D
diffusion is memory intensive, however. Wavelet diffusion models
(WDM)~\cite{friedrich2024wdm} reduce this cost through a fixed, invertible spatial
transform, and conditional WDM~\cite{friedrich2024cwdm} extends the formulation to
image-to-image tasks. We adapt it as CATCH (\textbf{C}ounterfactual
\textbf{A}natomical \textbf{T}issue Inpainting with \textbf{C}onditional
\textbf{H}aar-wavelet diffusion) and ask how healthy training holes should be
sampled.

Our contributions are:
\begin{itemize}
  \item conditional 3D Haar-wavelet diffusion with tumor-excluded supervision and
  voxel-exact preservation outside the requested hole;
  \item a within-CATCH comparison of fixed-mask, strict random, and
  weighted-mixture healthy-mask training variants, with compute-matched random
  versus weighted inference;
  \item development-only checkpoint and trajectory-policy selection followed by a
  frozen 75-case confirmation with paired uncertainty, patient-disjoint
  sensitivity, and audited examples.
\end{itemize}
